\chapter[只要将她铭记于心，你就并不孤单]{只要将她铭记于心，你就并不孤单}

\chaptersubtitle{《魔法少女小圆：The Different Story》}

颇为对称的是，《魔法少女小圆：The Different Story》（叶之影作画，原案 Magica Quartet）把重点放在外传作品必须处理的三项关键原则中最后尚未讨论的一项：
不公。
本作没有引入任何新角色，而是对动画采取一种范围颇为有限的“假如……会怎样”式处理。

第一卷主要刻画麻美与杏子，同时也是电视动画和漫画后续部分的前传。
在这一卷中，麻美收下刚刚成为魔法少女的杏子作为徒弟。
杏子的家人死后，两人的关系随之破裂，因为杏子不再认同麻美那套同时猎杀魔女与使魔、借此保护他人的理念。
第二卷与第三卷则转入一条动画未曾呈现的时间线，在这里，麻美收下沙耶香作为新的魔法少女徒弟。
由于沙耶香比动画中稍早与丘比订立契约，她得以及时救下麻美，使其免遭夏洛特杀害。
随后杏子归来，她与麻美、沙耶香围绕魔法少女应当以何种方式生存发生冲突；
与此同时，沙耶香逐渐陷入抑郁。

这部漫画反复审视魔法少女体制根本上的不公，但切入角度有别于《魔圆》。
电视动画主要关注丘比如何以不公正的方式向签约少女隐瞒关键信息，例如她们身体被改造的事实，或魔女的真正本质。
前一项在\emph{The Different Story}中完全没有出现，后一项虽然出现，却也只得到寥寥数语的讨论。
本作真正关注的，是强加在少女们身上的孤立，以及这种处境对她们精神状态的影响。
当然，孤立及其引发的抑郁，也是《魔圆》中沙耶香篇的重点。
但在动画中，沙耶香至少在一定程度上主动造成了自己的孤立：
得知身体遭到改造后，她有意将身边的人拒之门外。
而在\emph{The Different Story}中，孤立无从避免，是魔法少女体制内在的特征。

整部漫画里，麻美一直在设法摆脱自己的孤立处境。
她极度渴望友谊，但学生生活与魔法少女职责带来的无尽劳碌，使她没有时间同班级里的同学相处。
她试图主动接近其他魔法少女，却遭到拒绝，因为她“愚蠢地”与使魔战斗，而没有把全部精力放在狩猎魔女、获取随之而来的悲叹之种上。
她的确成功与杏子成为朋友，却又因为杏子选择了其他魔法少女共同奉行的那套理念，而失去了这位朋友。
后来，她又与小圆和沙耶香建立起感情，但沙耶香最终变成魔女，部分原因恰恰在于她遵循了麻美猎杀使魔的方针。
即便到了最后，小圆已经为麻美提供了一条出路——
代替她以魔法少女的身份战斗，让她得以休息——
麻美仍摆脱不了孤独与自我否定，最终走向自杀。

麻美的例子最为突出，但\emph{The Different Story}中另外两名逐渐陷入绝望的少女也如此。
家人死后，杏子几乎彻底封闭了自己，除了小圆，她拒绝所有试图与自己交朋友的人。
与此同时，沙耶香则因为曾把仁美弃于魔女之手而心怀羞耻，有意疏远了麻美、杏子与恭介。

这种孤立固然可以说是出于她们自己的选择，但漫画同时又将其归因于广义上的世界不公，尤其是丘比的魔法少女体制之不公。
贯穿整部漫画的麻美与其他魔法少女之间的核心争论，就在于获取悲叹之种是否必须凌驾于保护他人等其他考量之上。
不仅如此，丘比甚至明确表示，沙耶香之所以变成魔女，正是因为见泷原聚集了过多魔法少女。
因为在这种局面下，实力最弱、经验最浅的魔法少女必然得不到足够的悲叹之种。
麻美连一个二人小队都无法维系，并非偶然，而是这套体制自身的特征：
它的设计目的，就是最大限度地制造魔法少女的绝望，加速她们向魔女转化。
如果没有其他因素先让两名魔法少女分道扬镳，她们最终也必然会为悲叹之种相互争夺。

杏子的故事则把这种不公延伸到更广阔的世界。
她的父亲与丘比一样，将魔法少女和魔女等同起来，但又走得更远：
在他眼中，二者根本没有任何区别。
他一得知杏子变成了什么，便彻底否定了她，声称她欺骗了自己，也败坏了自己的教诲。
他只能把魔法或女性力量理解为巫术，并拒绝倾听杏子的解释。
随后，他开始酗酒，继而施暴，最终以杀人和自杀收场。

三位少女的经历都贯穿着一种错误的罪疚感，这一点强调了魔法少女并不应为自身的孤立承担责任。
麻美因家人的死亡而责备自己，因为她当时没有许愿拯救他们；
但一个濒死之人只想着活下去，并因此排除其他一切愿望，本来也很难说是不合理的。
与此相反，杏子因为家人的死亡而责备自己，原因恰恰在于她许愿拯救他们；
然而，她确实是在帮助他们，真正导致他们死亡的，只是她父亲拒绝倾听，也拒绝尝试理解。
这种罪疚感在经历悲伤、失落与创伤的人身上相当常见，也是创伤后应激障碍的一项主要症状\cite{apa2000dsm-iv-tr}。
沙耶香同样责备自己未能从魔女手中救下仁美；
但正如小圆指出的那样，她当时并不确定那个人就是仁美，也无从知道此事牵涉到魔女。

三名少女都因罪疚感而孤立自己。
沙耶香与杏子的表现更为明显：
她们会以言语攻击那些试图与自己成为朋友的人，有时甚至会动手伤害对方。
不过，麻美同样在孤立自己。
正如她在最终话中所揭示的那样，她创造出一个虚假的、看似“完美”的自我，避免让他人看见自己的需求。
而读者从杏子的视角中已经知道，这样做只会加剧他人的自卑感与孤立感，从而延续孤独与绝望的循环。
这种自我孤立也是创伤后应激障碍的另一项症状\cite{apa2000dsm-iv-tr}。
由此可见，漫画暗示丘比的系统本身就具有创伤性：
它使每一个进入这一系统的少女都承受战斗压力，同时遭受痛苦的孤立。

最终，小圆找到了一种解决办法：
她许愿让沙耶香复活，并取代麻美成为魔法少女。
但这仍然不够。
结尾的一连串场景看起来是麻美自杀后的临终之梦，而它表明，只有这样一个世界才足以抚慰麻美的孤独：
五名魔法少女——麻美、杏子、沙耶香、小圆和焰——
能够并肩合作、彼此为友，而麻美也拥有自己的家庭。

遗憾的是，这样的世界本身也并非没有问题。
这一点，我们将在讨论《叛逆》时看到。
